Dirac\cite{Dirac1936} has shown the existence of isomorphism between $SO(4,2)$ conformal group and Dirac matrices. Later, Hepner\cite{Hepner1962} has explicitly shown the isomorphism between the group of Dirac's four-row $\gamma$-matrices and the continuous conformal group in Euclidean space. One can show the isomorphism between the conformal group Eq.\eqref{Jab} and the group of fifteen matrices of the $\gamma$'s and their products in the Minkowski space,
{\begin{align}
  J'_{a,b}&=
  \scalebox{0.7}{\begin{pmatrix}
  0&\gamma_5&\frac{(1-\gamma_5)\gamma_0}{\sqrt{2}}&\frac{(1-\gamma_5)\gamma_1}{\sqrt{2}}&\frac{(1-\gamma_5)\gamma_2}{\sqrt{2}}&\frac{(1-\gamma_5)\gamma_3}{\sqrt{2}}\\
  -\gamma_5&0&\frac{(1+\gamma_5)\gamma_0}{\sqrt{2}}&\frac{(1+\gamma_5)\gamma_1}{\sqrt{2}}&\frac{(1+\gamma_5)\gamma_2}{\sqrt{2}}&\frac{(1+\gamma_5)\gamma_3}{\sqrt{2}}\\
    \frac{-(1-\gamma_5)\gamma_0}{\sqrt{2}}&\frac{-(1+\gamma_5)\gamma_0}{\sqrt{2}}&0 & \gamma_0\gamma_1 & \gamma_0\gamma_2 & \gamma_0\gamma_3\\
   \frac{-(1-\gamma_5)\gamma_1}{\sqrt{2}}&\frac{-(1+\gamma_5)\gamma_1}{\sqrt{2}}&\gamma_1\gamma_0 & 0 & \gamma_1\gamma_2 & \gamma_1\gamma_3\\
    \frac{-(1-\gamma_5)\gamma_2}{\sqrt{2}}&\frac{-(1+\gamma_5)\gamma_2}{\sqrt{2}}&\gamma_2\gamma_0 & \gamma_2\gamma_1 & 0 & \gamma_2\gamma_3\\
    \frac{-(1-\gamma_5)\gamma_3}{\sqrt{2}}&\frac{-(1+\gamma_5)\gamma_3}{\sqrt{2}}&\gamma_3\gamma_0 & \gamma_3\gamma_1 & \gamma_3\gamma_2 & 0
  \end{pmatrix}}. 
\end{align}}
This $J'_{a,b}$ obeys the $SO(4+1,1)$ algebra Eq.\eqref{algebrasimp} with the metric is given by Eq.\eqref{metric}. The representation of the conformal group in terms of $4\times4$ gamma matrices are the following;
\begin{align}
    P_\mu=&\frac{i}{2}(1+\gamma_5)\gamma_\mu;\\
    \mathfrak{K}_\mu=&\frac{-i}{2}(1-\gamma_5)\gamma_\mu;\\
    K^1=&\frac{i}{2}\gamma_1\gamma_0;~~K^2=\frac{i}{2}\gamma_2\gamma_0;~~K^3=\frac{i}{2}\gamma_3\gamma_0;\\
    J^1=&\frac{i}{2}\gamma_2\gamma_3;~~J^2=\frac{i}{2}\gamma_3\gamma_1;~~J^3=\frac{i}{2}\gamma_1\gamma_2;\\
    D=&\frac{-i}{2}\gamma_5.
\end{align}

\subsection{Isomorphism with Dirac matrices}
The defining property for the gamma matrices to generate a Clifford algebra is the anticommutation relation 
\begin{align}
    \left\{\gamma ^{\mu },\gamma ^{\nu }\right\}=\gamma ^{\mu }\gamma ^{\nu }+\gamma ^{\nu }\gamma ^{\mu }=2g ^{\mu \nu }I_{4},
\end{align}
Covariant gamma matrices are defined by
\begin{align}
    {\gamma _{\mu }=\eta _{\mu \nu }\gamma ^{\nu }=\left\{\gamma ^{0},-\gamma ^{1},-\gamma ^{2},-\gamma ^{3}\right\}.}
\end{align}
Further basis $\sigma _{\mu \nu }$ elements of the Clifford algebra are given by
\begin{align}
    \sigma _{\mu \nu }&={\frac {1}{4}}[\gamma _{\mu },\gamma _{\nu }]=-{\frac {i}{4}}(\gamma _{\mu }\gamma _{\nu }-\gamma _{\nu }\gamma _{\mu })\\
    \sigma _{\mu \nu }&={\frac {1}{4}}\left(\gamma _{\mu }\gamma _{\nu }-2g_{\mu\nu}+\gamma _{\mu }\gamma _{\nu }\right)\\
    \sigma _{\mu \nu }&={\frac {1}{2}}(\gamma _{\mu }\gamma _{\nu }-g_{\mu\nu})
\end{align}
where it satisfy,
\begin{align}
   \sigma_{\mu\nu}=&-\sigma_{\nu\mu};\\
   \left[\sigma _{\mu \nu },\sigma _{\rho \sigma }\right]=&\left(-g_{{\nu}{\sigma}}\sigma_{{\mu}{\rho}}+g_{{\nu}{\rho}}\sigma_{{\mu}{\sigma}}-g_{{\mu}{\rho}}\sigma_{{\nu}{\sigma}}+g_{{\mu}{\sigma}}\sigma_{{\nu}{\rho}}\right).
\end{align}

Let's go to $d=6$; we have $\sigma_{a,b}$ where $a,b\in\{-2,-1,0,1,2,3\}$.

Then $\gamma_a=(~~,\gamma_5,\gamma_0,\gamma_1,\gamma_2,\gamma_3)$.

Let's define 
\begin{align}
    J_{a,b}\equiv i\sigma_{a,b}
\end{align}
then $\left[\sigma _{ab},\sigma _{cd}\right]=&\left(-g_{{b}{d}}\sigma_{{a}{c}}+g_{{b}{c}}\sigma_{{a}{d}}-g_{{a}{c}}\sigma_{{b}{d}}+g_{{a}{d}}\sigma_{{b}{c}}\right)$ will be (since, $\sigma_{a,b}=-iJ_{a,b}$)
\begin{align}
    (-1)\left[J _{ab},J _{cd}\right]=&(-i)\left(-g_{{b}{d}}J_{{a}{c}}+g_{{b}{c}}J_{{a}{d}}-g_{{a}{c}}J_{{b}{d}}+g_{{a}{d}}J_{{b}{c}}\right)\\
    \left[J _{ab},J _{cd}\right]=&-i\left(g_{{b}{d}}J_{{a}{c}}-g_{{b}{c}}J_{{a}{d}}+g_{{a}{c}}J_{{b}{d}}-g_{{a}{d}}J_{{b}{c}}\right)
\end{align}
\begin{align}
  J_{a,b}&={\frac {i}{2}}(\gamma _{\mu }\gamma _{\nu }-g_{\mu\nu})=\frac{i}{2}
  \begin{pmatrix}
  0&\gamma_5&\gamma_0&\gamma_1&\gamma_2&\gamma_3\\
  -\gamma_5&0&\gamma_5\gamma_0&\gamma_5\gamma_1&\gamma_5\gamma_2&\gamma_5\gamma_3\\
    -\gamma_0&\gamma_0\gamma_5&0 & \gamma_0\gamma_1 & \gamma_0\gamma_2 & \gamma_0\gamma_3\\
    -\gamma_1&\gamma_1\gamma_5&\gamma_1\gamma_0 & 0 & \gamma_1\gamma_2 & \gamma_1\gamma_3\\
    -\gamma_2&\gamma_2\gamma_5&\gamma_2\gamma_0 & \gamma_2\gamma_1 & 0 & \gamma_2\gamma_3\\
    -\gamma_3&\gamma_3\gamma_5&\gamma_3\gamma_0 & \gamma_3\gamma_1 & \gamma_3\gamma_2 & 0
  \end{pmatrix}_{6\times6}
\end{align}
  \begin{align}
      g_{ab}=\begin{pmatrix}
  -1&0&0&0&0&0\\
  0&1&0&0&0&0\\
  0&0&1&0&0&0\\
  0&0&0&-1&0&0\\
  0&0&0&0&-1&0\\
  0&0&0&0&0&-1\\
  \end{pmatrix}_{6\times6}
  \end{align}

After the 6 dimensional rotation,

\begin{align}
  J'_{a,b}&=
  \begin{pmatrix}
  0&D&\frac{-\mathfrak{K}_0}{\sqrt{2}}&\frac{-\mathfrak{K}_1}{\sqrt{2}}&\frac{-\mathfrak{K}_2}{\sqrt{2}}&\frac{-\mathfrak{K}_3}{\sqrt{2}}\\
  -D&0&\frac{P_0}{\sqrt{2}}&\frac{P_1}{\sqrt{2}}&\frac{P_2}{\sqrt{2}}&\frac{P_3}{\sqrt{2}}\\
    \frac{\mathfrak{K}_0}{\sqrt{2}}&\frac{-P_0}{\sqrt{2}}&0 & -K^{1} & -K^{2} & -K^{3}\\
    \frac{\mathfrak{K}_1}{\sqrt{2}}&\frac{-P_1}{\sqrt{2}}&K^{1} & 0 & J^{3} & -J^{2}\\
    \frac{\mathfrak{K}_2}{\sqrt{2}}&\frac{-P_2}{\sqrt{2}}&K^{2} & -J^{3} & 0 & J^{1}\\
    \frac{\mathfrak{K}_3}{\sqrt{2}}&\frac{-P_3}{\sqrt{2}}&K^{3} & J^{2} & -J^{1} & 0
  \end{pmatrix}_{6\times6}
\end{align}

\begin{align}
  \boxed{J'_{a,b}=\frac{i}{2}
  \begin{pmatrix}
  0&\gamma_5&\frac{(1-\gamma_5)\gamma_0}{\sqrt{2}}&\frac{(1-\gamma_5)\gamma_1}{\sqrt{2}}&\frac{(1-\gamma_5)\gamma_2}{\sqrt{2}}&\frac{(1-\gamma_5)\gamma_3}{\sqrt{2}}\\
  -\gamma_5&0&\frac{(1+\gamma_5)\gamma_0}{\sqrt{2}}&\frac{(1+\gamma_5)\gamma_1}{\sqrt{2}}&\frac{(1+\gamma_5)\gamma_2}{\sqrt{2}}&\frac{(1+\gamma_5)\gamma_3}{\sqrt{2}}\\
    \frac{-(1-\gamma_5)\gamma_0}{\sqrt{2}}&\frac{-(1+\gamma_5)\gamma_0}{\sqrt{2}}&0 & \gamma_0\gamma_1 & \gamma_0\gamma_2 & \gamma_0\gamma_3\\
   \frac{-(1-\gamma_5)\gamma_1}{\sqrt{2}}&\frac{-(1+\gamma_5)\gamma_1}{\sqrt{2}}&\gamma_1\gamma_0 & 0 & \gamma_1\gamma_2 & \gamma_1\gamma_3\\
    \frac{-(1-\gamma_5)\gamma_2}{\sqrt{2}}&\frac{-(1+\gamma_5)\gamma_2}{\sqrt{2}}&\gamma_2\gamma_0 & \gamma_2\gamma_1 & 0 & \gamma_2\gamma_3\\
    \frac{-(1-\gamma_5)\gamma_3}{\sqrt{2}}&\frac{-(1+\gamma_5)\gamma_3}{\sqrt{2}}&\gamma_3\gamma_0 & \gamma_3\gamma_1 & \gamma_3\gamma_2 & 0
  \end{pmatrix}_{6\times6}}
\end{align}

  \begin{align}
      g_{a,b}=\begin{pmatrix}
  0&-1&0&0&0&0\\
  -1&0&0&0&0&0\\
  0&0&1&0&0&0\\
  0&0&0&-1&0&0\\
  0&0&0&0&-1&0\\
  0&0&0&0&0&-1\\
  \end{pmatrix}_{6\times6}
  \end{align}
So fifteen matrices of the $\gamma$'s and their products can be grouped as the fifteen components of a skew angular momentum tensor
$J_{a,b}$ in six dimensions\footnote{The inhomogeneous lorentz group and the conformal group,
W. A. Hepner}.

\\

The representation of conformal group in terms of $4\times4$ gamma matrices are the following;
\begin{align}
    P_\mu=&\frac{i}{2}(1+\gamma_5)\gamma_\mu;\\
    \mathfrak{K}_\mu=&\frac{-i}{2}(1-\gamma_5)\gamma_\mu;\\
    K^1=&\frac{i}{2}\gamma_1\gamma_0;~~K^2=\frac{i}{2}\gamma_2\gamma_0;~~K^3=\frac{i}{2}\gamma_3\gamma_0;\\
    J^1=&\frac{i}{2}\gamma_2\gamma_3;~~J^2=\frac{i}{2}\gamma_3\gamma_1;~~J^3=\frac{i}{2}\gamma_1\gamma_2;\\
    D=&\frac{i}{2}\gamma_5.
\end{align}
Explicitly,
\begin{align*}
    J^{1}&=\frac{1}{2}\begin{pmatrix}
    0&1&0&0\\
    1&0&0&0\\
    0&0&0&1\\
    0&0&1&0
    \end{pmatrix}~;~~~~~~J^{2}=\frac{1}{2}\begin{pmatrix}
    0&-i&0&0\\
    i&0&0&0\\
    0&0&0&-i\\
    0&0&i&0
    \end{pmatrix}~;~~~~~~J^{3}=\frac{1}{2}\begin{pmatrix}
    1&0&0&0\\
    0&-1&0&0\\
    0&0&1&0\\
    0&0&0&-1
    \end{pmatrix};\\
    K^{1}&=\frac{1}{2}\begin{pmatrix}
    0&i&0&0\\
    i&0&0&0\\
    0&0&0&-i\\
    0&0&-i&0
    \end{pmatrix}~;~~~~~~K^{2}=\frac{1}{2}\begin{pmatrix}
    0&1&0&0\\
    -1&0&0&0\\
    0&0&0&-1\\
    0&0&1&0
    \end{pmatrix}~;~~~~~~K^{3}=\frac{1}{2}\begin{pmatrix}
    i&0&0&0\\
    0&-i&0&0\\
    0&0&-i&0\\
    0&0&0&i
    \end{pmatrix};\\
    P_0&=\begin{pmatrix}
    0&0&0&0\\
    0&0&0&0\\
    i&0&0&0\\
    0&i&0&0
    \end{pmatrix}~;~~~~~~P_1=\begin{pmatrix}
    0&0&0&0\\
    0&0&0&0\\
    0&-i&0&0\\
    -i&0&0&0
    \end{pmatrix}~;~~~~~~P_2=\begin{pmatrix}
    0&0&0&0\\
    0&0&0&0\\
    0&-1&0&0\\
    1&0&0&0
    \end{pmatrix}~;~~~~~~P_3=\begin{pmatrix}
    0&0&0&0\\
    0&0&0&0\\
    -i&0&0&0\\
    0&i&0&0
    \end{pmatrix}\\
    \mathfrak{K}_0&=\begin{pmatrix}
    0&0&-i&0\\
    0&0&0&-i\\
    0&0&0&0\\
    0&0&0&0
    \end{pmatrix}~;~~~~~~\mathfrak{K}_1=\begin{pmatrix}
    0&0&0&-i\\
    0&0&-i&0\\
    0&0&0&0\\
    0&0&0&0
    \end{pmatrix}~;~~~~~~\mathfrak{K}_2=\begin{pmatrix}
    0&0&0&-1\\
    0&0&1&0\\
    0&0&0&0\\
    0&0&0&0
    \end{pmatrix}~;~~~~~~\mathfrak{K}_3=\begin{pmatrix}
    0&0&-i&0\\
    0&0&0&i\\
    0&0&0&0\\
    0&0&0&0
    \end{pmatrix}\\
    D&=\frac{i}{2}\begin{pmatrix}
    -1&0&0&0\\
    0&-1&0&0\\
    0&0&1&0\\
    0&0&0&1
    \end{pmatrix}.
\end{align*}

this is dirac paper:
@article{436861d5-35f6-36cf-9470-cb587eced490,
 ISSN = {0003486X, 19398980},
 URL = {http://www.jstor.org/stable/1968455},
 author = {P. A. M. Dirac},
 journal = {Annals of Mathematics},
 number = {2},
 pages = {429--442},
 publisher = {[Annals of Mathematics, Trustees of Princeton University on Behalf of the Annals of Mathematics, Mathematics Department, Princeton University]},
 title = {Wave Equations in Conformal Space},
 urldate = {2026-08-26},
 volume = {37},
 year = {1936}
}

this is hepner:
@article{Hepner1962The,
  author = {Hepner, W. A.},
  title = {The inhomogeneous lorentz group and the conformal group},
  journal = {Il Nuovo Cimento (1955-1965)},
  year = {1962},
  volume = {26},
  number = {2},
  pages = {351--368},
  doi = {10.1007/BF02787046},
  url = {https://doi.org/10.1007/BF02787046},
  isbn = {1827-6121}
}

